\documentclass[12pt,a4paper]{article}

\usepackage[margin=2.5cm]{geometry}
\usepackage{amsmath, amssymb}
\usepackage{booktabs}
\usepackage{graphicx}
\usepackage{hyperref}
\usepackage{listings}
\usepackage{xcolor}
\usepackage{float}
\usepackage{caption}
\usepackage{subcaption}
\usepackage{parskip}
\usepackage{tikz}
\usetikzlibrary{shapes.geometric, arrows.meta, positioning, fit, backgrounds}

\hypersetup{colorlinks=true, linkcolor=blue, urlcolor=blue}

\lstset{
  basicstyle=\ttfamily\small,
  backgroundcolor=\color{gray!10},
  frame=single,
  breaklines=true,
  keywordstyle=\color{blue},
  commentstyle=\color{green!50!black},
  stringstyle=\color{red!60!black},
}

\title{\textbf{Model 4: Genre \& Keywords Bayesian Regression} \\
\large Probabilistic Movie Revenue Prediction via ARD Priors and SVI}
\author{Arham Aziz Noman}
\date{May 2026}

\begin{document}

\maketitle
\tableofcontents
\newpage

% ─────────────────────────────────────────────────────────────────────────────
\section{Overview}

Model 4 predicts movie box-office revenue using two content-based signals:
\textbf{genre} (what category of film it is) and \textbf{keywords} (free-text tags
describing its content). It is a Bayesian linear regression model with
Automatic Relevance Determination (ARD) priors, trained using Stochastic
Variational Inference (SVI) in Pyro.

The key output is not a single number but a full probability distribution
$\mathcal{N}(\mu_{\text{genre}},\, \sigma_{\text{genre}}^2)$ per movie.
This distribution is consumed by a precision-weighted Bayesian aggregator
that fuses predictions from multiple models.

% ─────────────────────────────────────────────────────────────────────────────
\section{Data, Features, and Model}

\subsection{Model Pipeline Diagram}

\begin{figure}[H]
\centering
\begin{tikzpicture}[
  node distance = 0.55cm and 1.1cm,
  every node/.style = {font=\small},
  % styles
  data/.style   = {draw, rounded corners=3pt, fill=blue!10,  minimum width=3.2cm, minimum height=0.7cm, align=center},
  feat/.style   = {draw, rounded corners=3pt, fill=green!12, minimum width=3.2cm, minimum height=0.7cm, align=center},
  model/.style  = {draw, rounded corners=3pt, fill=orange!15,minimum width=3.6cm, minimum height=0.7cm, align=center},
  prior/.style  = {draw, rounded corners=3pt, fill=red!10,   minimum width=3.2cm, minimum height=0.7cm, align=center},
  out/.style    = {draw, rounded corners=3pt, fill=purple!12,minimum width=3.2cm, minimum height=0.7cm, align=center},
  arr/.style    = {-{Stealth[length=5pt]}, thick},
  grp/.style    = {draw=gray!50, dashed, rounded corners=5pt, inner sep=6pt},
]

%% ── Column 1: raw inputs ──────────────────────────────────────────
\node[data] (meta)  {movies\_metadata\_slim.csv\\{\tiny 5369 rows, 7 cols}};
\node[data, below=of meta] (kw)    {keywords.csv\\{\tiny 46k movies, tags}};
\node[data, below=of kw]   (cpi)   {US CPI 1960--2024\\{\tiny inflation rates}};

%% ── Column 2: preprocessing ──────────────────────────────────────
\node[feat, right=of meta]  (infl)  {CPI Deflation\\{\tiny revenue $\to$ 2024 USD}};
\node[feat, below=of infl]  (log)   {Log Transform\\{\tiny $y=\log(\text{rev})$}};
\node[feat, below=of log]   (tfidf) {TF-IDF\\{\tiny 500 keywords, min\_df=3}};
\node[feat, below=of tfidf] (svd)   {Truncated SVD\\{\tiny $500 \to 20$ topics}};

%% ── Column 3: feature groups ─────────────────────────────────────
\node[feat, right=of infl,  yshift=-0.6cm] (fgenre)  {Genre one-hot\\{\tiny $\mathbf{x}_g \in \mathbb{R}^{15}$, weighted}};
\node[feat, below=of fgenre]               (fkw)     {Keyword topics\\{\tiny $\mathbf{x}_k \in \mathbb{R}^{20}$, scaled}};
\node[feat, below=of fkw]                  (frt)     {Log-runtime\\{\tiny $x_r \in \mathbb{R}^{1}$, scaled}};

%% ── Column 4: Bayesian model ──────────────────────────────────────
\node[model, right=of fgenre, yshift=-0.6cm] (ard1) {$\sigma_g \sim \text{HN}(5)$\\$\boldsymbol{\beta}_g \sim \mathcal{N}(\mathbf{0},\sigma_g^2\mathbf{I})$};
\node[model, below=of ard1]                  (ard2) {$\sigma_k \sim \text{HN}(2)$\\$\boldsymbol{\beta}_k \sim \mathcal{N}(\mathbf{0},\sigma_k^2\mathbf{I})$};
\node[model, below=of ard2]                  (brt)  {$\beta_r \sim \mathcal{N}(0,4)$};

%% ── Linear predictor ──────────────────────────────────────────────
\node[model, right=of ard1, yshift=-0.9cm, minimum width=3.8cm] (mu)
  {$\mu_i = \alpha + \mathbf{x}_g^\top\boldsymbol{\beta}_g + \mathbf{x}_k^\top\boldsymbol{\beta}_k + x_r\beta_r$};

%% ── Likelihood ────────────────────────────────────────────────────
\node[prior, below=of mu] (lik) {$y_i \sim \mathcal{N}(\mu_i,\,\sigma_\text{obs}^2)$\\{\tiny $\sigma_\text{obs}\sim\text{HN}(1)$}};

%% ── Output ────────────────────────────────────────────────────────
\node[out, below=of lik] (out) {Output: $\mathcal{N}(\hat{\mu},\,\hat{\sigma}^2)$\\{\tiny per movie — 500 posterior samples}};
\node[out, below=of out] (prec) {Precision: $1/\hat{\sigma}^2$\\{\tiny weight in Bayesian aggregator}};

%% ── Arrows: inputs → preprocessing ───────────────────────────────
\draw[arr] (meta) -- (infl);
\draw[arr] (cpi)  -- (infl);
\draw[arr] (meta) -- (log);
\draw[arr] (kw)   -- (tfidf);
\draw[arr] (tfidf)-- (svd);

%% ── Arrows: preprocessing → feature groups ────────────────────────
\draw[arr] (meta)  -- (fgenre);
\draw[arr] (log)   -- (fgenre);
\draw[arr] (svd)   -- (fkw);
\draw[arr] (meta)  -| (frt);

%% ── Arrows: features → ARD priors ────────────────────────────────
\draw[arr] (fgenre) -- (ard1);
\draw[arr] (fkw)    -- (ard2);
\draw[arr] (frt)    -- (brt);

%% ── Arrows: ARD → linear predictor ──────────────────────────────
\draw[arr] (ard1) -- (mu);
\draw[arr] (ard2) -- (mu);
\draw[arr] (brt)  -- (mu);

%% ── Arrows: predictor → output ───────────────────────────────────
\draw[arr] (mu)  -- (lik);
\draw[arr] (lik) -- (out);
\draw[arr] (out) -- (prec);

%% ── Background group boxes ───────────────────────────────────────
\begin{scope}[on background layer]
  \node[grp, fit=(meta)(kw)(cpi), label=above:{\textbf{Raw Data}}] {};
  \node[grp, fit=(infl)(log)(tfidf)(svd), label=above:{\textbf{Preprocessing}}] {};
  \node[grp, fit=(fgenre)(fkw)(frt), label=above:{\textbf{Features}}] {};
  \node[grp, fit=(ard1)(ard2)(brt)(mu)(lik), label=above:{\textbf{Bayesian Model}}] {};
  \node[grp, fit=(out)(prec), label=above:{\textbf{Output}}] {};
\end{scope}

\end{tikzpicture}
\caption{Model 4 pipeline: from raw data to posterior predictive output}
\end{figure}

\subsection{Dataset}

The model uses the TMDB Full Movie Dataset. Two files are required:

\begin{table}[H]
\centering
\begin{tabular}{lll}
\toprule
\textbf{File} & \textbf{Size} & \textbf{Contents} \\
\midrule
\texttt{movies\_metadata\_slim.csv} & 0.8 MB & 5,369 rows, 7 columns \\
\texttt{keywords.csv}               & 6 MB   & 46,419 movies, keyword tags \\
\texttt{US CPI 1960--2024.csv}      & small  & Annual inflation rates \\
\bottomrule
\end{tabular}
\caption{Input data files}
\end{table}

\texttt{movies\_metadata\_slim.csv} is a pre-filtered version of the full 34 MB
metadata file, retaining only the columns needed by this model:
\texttt{id}, \texttt{title}, \texttt{budget}, \texttt{revenue},
\texttt{runtime}, \texttt{genres}, \texttt{release\_date}.
Rows with missing or zero budget, revenue, or runtime are excluded at
creation time, resulting in 5,369 usable movies.

\subsection{Inflation Adjustment}

All monetary values are expressed in \textbf{2024 USD} using US Consumer
Price Index (CPI) data. Without this adjustment, a film earning \$100M in
1990 and one earning \$100M in 2020 would appear equally valuable to the
model, despite \$100M in 1990 being worth approximately \$230M today.

The deflator for year $t$ is computed as the cumulative product of annual
inflation factors from year $t$ to 2024:

\begin{equation}
  \text{deflator}(t) = \prod_{y=t}^{2024} \bigl(1 + r_y\bigr)
\end{equation}

where $r_y$ is the inflation rate in year $y$. Adjusted revenue is then:

\begin{equation}
  \text{revenue\_2024}(t) = \text{revenue}(t) \times \text{deflator}(t)
\end{equation}

\subsection{Log Transformation}

Revenue and budget span several orders of magnitude (\$thousands to
\$billions). A linear model applied to raw values would be dominated by
a handful of blockbusters. Log-transforming compresses the scale:

\begin{equation}
  y = \log\bigl(\text{revenue\_2024}\bigr)
\end{equation}

This also makes the Normal likelihood assumption appropriate:
$\log(\text{revenue})$ is approximately bell-shaped, whereas raw revenue
is heavily right-skewed. After filtering and log-transformation, the
target $y$ ranges from 0.28 to 22.18 (log-scale).

\subsection{Train / Test Split}

\begin{table}[H]
\centering
\begin{tabular}{lrr}
\toprule
\textbf{Split} & \textbf{Movies} & \textbf{Fraction} \\
\midrule
Training & 4,152 & 80\% \\
Test     & 1,038 & 20\% \\
\midrule
Total    & 5,190 & 100\% \\
\bottomrule
\end{tabular}
\caption{Train/test split (random seed = 42). Note: 5,190 rows pass the
inflation-adjustment filter; 179 rows are dropped due to missing CPI years.}
\end{table}

\subsection{Genre Features (One-Hot Weighted)}

Each movie is assigned to one or more genres by TMDB. The top 15 most
frequent genres are retained. A multi-genre film is represented as a
\emph{distribution} over genres rather than a binary flag vector: each
genre indicator is divided by the total number of genres assigned to that film.

\begin{equation}
  x_g^{(i)} =
  \begin{cases}
    \dfrac{1}{|\mathcal{G}_i|} & \text{if genre } g \in \mathcal{G}_i \\[6pt]
    0                          & \text{otherwise}
  \end{cases}
\end{equation}

where $\mathcal{G}_i$ is the set of genres assigned to movie $i$.
This normalization ensures $\sum_g x_g^{(i)} = 1$ for every movie,
preventing multi-genre films from having higher total feature magnitude
than single-genre films.

\noindent\textbf{Top 15 genres (by frequency):}
Drama, Comedy, Thriller, Action, Romance, Adventure, Crime,
Science Fiction, Horror, Family, Fantasy, Mystery, Animation, History, War.

\subsection{Keyword Topic Features (TF-IDF $\rightarrow$ SVD)}

TMDB provides free-text keyword tags per movie (e.g., \textit{based\_on\_novel},
\textit{sequel}, \textit{superhero}). Keywords capture content signals that
go beyond genre labels.

\paragraph{Step 1 — TF-IDF.}
Each movie's keywords are concatenated into a single document.
TF-IDF weighting is applied over a vocabulary of the 500 most frequent
keywords that appear in at least 3 movies. Keywords appearing in nearly
all movies receive low weight; distinctive keywords receive high weight.

\paragraph{Step 2 — Truncated SVD (LSA).}
The $5190 \times 500$ sparse TF-IDF matrix is compressed to a dense
$5190 \times 20$ matrix using Truncated SVD (Latent Semantic Analysis).
Each of the 20 latent topics captures a cluster of co-occurring keywords.
The 20 topics explain 18.56\% of variance in keyword co-occurrence patterns.

\paragraph{Step 3 — Standardization.}
SVD components can have very different magnitudes across topics.
StandardScaler (zero-mean, unit-variance) is applied to put all topics
on the same scale.

\subsection{Runtime Feature}

Log-runtime is computed as $\log(\text{runtime\_minutes})$ and then
standardized (zero-mean, unit-variance). The log-transform is used because
runtime spans a large range and its relationship with revenue is likely
multiplicative rather than additive.

\subsection{Feature Summary}

\begin{table}[H]
\centering
\begin{tabular}{llrr}
\toprule
\textbf{Feature group} & \textbf{Transformation} & \textbf{Dim} & \textbf{Scaling} \\
\midrule
Genre        & Weighted one-hot (top 15) & 15 & None (already in $[0,1]$) \\
Keywords     & TF-IDF $\to$ SVD topics   & 20 & StandardScaler \\
Runtime      & Log-minutes               &  1 & StandardScaler \\
\midrule
Total        & & 36 & \\
\bottomrule
\end{tabular}
\caption{Feature matrix: $N \times 36$ total input to the model}
\end{table}

\subsection{Why Bayesian Regression?}

A standard linear regression returns a single point estimate.
A Bayesian model returns a full posterior distribution over the weights,
and therefore a full \emph{predictive distribution} $p(y^* \mid x^*)$ for
any new input $x^*$. The predictive uncertainty $\sigma$ is the central
output of Model 4: it is used as the inverse-variance weight in the
Bayesian aggregator that fuses all models.

\subsection{Why ARD Priors?}

Automatic Relevance Determination (ARD) assigns each feature \emph{group}
its own shared scale hyperparameter. If the data provides no evidence that
genre predicts revenue, the posterior $\sigma_\text{genre}$ stays small
and all genre weights shrink toward zero. If genre is informative,
$\sigma_\text{genre}$ grows and allows large $\beta$ weights.
The model thus learns how much to trust each feature group automatically,
without manual regularization tuning.

\subsection{Generative Model}

\begin{align}
  \alpha &\sim \mathcal{N}(14,\; 3^2)
    && \text{intercept (log-revenue $\approx 14$--$20$)} \\[4pt]
  \sigma_\text{genre} &\sim \text{HalfNormal}(5)
    && \text{ARD scale for genre group} \\
  \boldsymbol{\beta}_\text{genre} &\sim \mathcal{N}\!\left(
      \mathbf{0},\; \sigma_\text{genre}^2\,\mathbf{I}_{15}
    \right)
    && \text{15 genre weights} \\[4pt]
  \sigma_\text{kw} &\sim \text{HalfNormal}(2)
    && \text{ARD scale for keyword group} \\
  \boldsymbol{\beta}_\text{kw} &\sim \mathcal{N}\!\left(
      \mathbf{0},\; \sigma_\text{kw}^2\,\mathbf{I}_{20}
    \right)
    && \text{20 keyword weights} \\[4pt]
  \beta_\text{rt} &\sim \mathcal{N}(0,\; 2^2)
    && \text{runtime weight} \\[4pt]
  \sigma_\text{obs} &\sim \text{HalfNormal}(1)
    && \text{observation noise} \\[8pt]
  \mu_i &= \alpha
    + \mathbf{x}_{\text{genre}}^{(i)\top}\boldsymbol{\beta}_\text{genre}
    + \mathbf{x}_{\text{kw}}^{(i)\top}\boldsymbol{\beta}_\text{kw}
    + x_{\text{rt}}^{(i)}\,\beta_\text{rt} \\[4pt]
  y_i &\sim \mathcal{N}(\mu_i,\; \sigma_\text{obs}^2)
    && \text{log-revenue likelihood}
\end{align}

\subsection{Prior Justification}

\begin{table}[H]
\centering
\begin{tabular}{lll}
\toprule
\textbf{Parameter} & \textbf{Prior} & \textbf{Reason} \\
\midrule
$\alpha$ & $\mathcal{N}(14, 9)$ & Log-revenue ranges 0--22; centred at typical value \\
$\sigma_\text{genre}$ & HalfNormal(5) & Wide — genre can have large effects \\
$\sigma_\text{kw}$ & HalfNormal(2) & Tighter — keyword topics are noisier signals \\
$\beta_\text{rt}$ & $\mathcal{N}(0, 4)$ & Weakly informative fixed prior for runtime \\
$\sigma_\text{obs}$ & HalfNormal(1) & Residual noise on log-scale \\
\bottomrule
\end{tabular}
\caption{Prior distributions and justifications}
\end{table}

% ─────────────────────────────────────────────────────────────────────────────
\section{Inference: Stochastic Variational Inference}

\subsection{Why SVI instead of MCMC?}

Markov Chain Monte Carlo (MCMC) gives exact samples from the posterior
but is computationally prohibitive for $N = 4{,}152$ training examples
and 36 parameters. Stochastic Variational Inference (SVI) approximates
the posterior with a tractable variational family and optimises it using
gradient descent — orders of magnitude faster while still providing
calibrated uncertainty estimates.

\subsection{Variational Family}

A \textbf{mean-field Gaussian} (AutoNormal) is used: each latent variable
$\theta_k$ is approximated independently as:

\begin{equation}
  q(\theta_k) = \mathcal{N}(\mu_k,\; \sigma_k^2)
\end{equation}

This introduces one variational mean $\mu_k$ and one variational
standard deviation $\sigma_k$ per latent variable.

\subsection{ELBO Objective}

SVI maximises the Evidence Lower BOund (ELBO):

\begin{equation}
  \mathcal{L}(q) = \mathbb{E}_{q(\boldsymbol{\theta})}
    \bigl[\log p(\mathbf{y} \mid \mathbf{X}, \boldsymbol{\theta})\bigr]
    - \mathrm{KL}\bigl(q(\boldsymbol{\theta}) \,\|\, p(\boldsymbol{\theta})\bigr)
\end{equation}

Maximising the ELBO simultaneously:
\begin{enumerate}
  \item Fits the data well (maximises the expected log-likelihood)
  \item Keeps the posterior close to the prior (minimises KL divergence)
\end{enumerate}

\subsection{Training Setup}

\begin{table}[H]
\centering
\begin{tabular}{ll}
\toprule
\textbf{Hyperparameter} & \textbf{Value} \\
\midrule
Optimizer       & ClippedAdam (gradient clip norm = 5.0) \\
Learning rate   & 0.01 \\
SVI steps       & 3,000 \\
Batch size      & Full dataset (no mini-batching) \\
Random seed     & 42 \\
\bottomrule
\end{tabular}
\caption{Training hyperparameters}
\end{table}

The ELBO loss converges from $\approx663{,}000$ at step 0 to
$\approx9{,}560$ at step 2,500, stabilising well before 3,000 steps.

\begin{figure}[H]
  \centering
  \fbox{\parbox{0.6\textwidth}{\centering\vspace{2cm}
    \textit{[Insert SVI convergence plot here]} \\
    ELBO loss vs. training step
  \vspace{2cm}}}
  \caption{SVI convergence: ELBO loss over 3,000 training steps}
\end{figure}

% ─────────────────────────────────────────────────────────────────────────────
\section{Posterior Predictive}

After training, the posterior predictive distribution for a new movie $x^*$
is approximated by drawing $S = 500$ samples from the variational posterior:

\begin{equation}
  p(y^* \mid x^*) \approx \frac{1}{S} \sum_{s=1}^{S}
    \mathcal{N}\!\left(y^* \;\big|\; \mu(x^*, \boldsymbol{\theta}^{(s)}),\;
    \sigma_\text{obs}^{(s)\,2}\right), \quad
    \boldsymbol{\theta}^{(s)} \sim q(\boldsymbol{\theta})
\end{equation}

The mean and standard deviation across the $S$ samples give the model output:

\begin{equation}
  \hat{\mu}_i = \frac{1}{S}\sum_s y_i^{(s)}, \qquad
  \hat{\sigma}_i = \sqrt{\frac{1}{S}\sum_s \bigl(y_i^{(s)} - \hat{\mu}_i\bigr)^2}
\end{equation}

% ─────────────────────────────────────────────────────────────────────────────
\section{Results}

\subsection{Predictive Accuracy}

\begin{table}[H]
\centering
\begin{tabular}{lcc}
\toprule
\textbf{Metric} & \textbf{Train} & \textbf{Test} \\
\midrule
RMSE (log-revenue)           & 2.3558 & 2.2225 \\
Mean log-likelihood (test)   & ---    & $-2.2206$ \\
\bottomrule
\end{tabular}
\caption{Model performance. RMSE is on the log-revenue scale.
An RMSE of 2.22 on log-scale corresponds to roughly a factor of
$e^{2.22} \approx 9\times$ error in revenue space.}
\end{table}

\noindent The test RMSE being lower than training RMSE suggests mild
underfitting — the ARD priors are regularising the model, and the
variational approximation introduces some bias. The model correctly
captures the broad revenue structure from genre and keywords but cannot
account for star power, marketing spend, or release timing.

\begin{figure}[H]
  \centering
  \fbox{\parbox{0.6\textwidth}{\centering\vspace{2cm}
    \textit{[Insert predicted vs actual scatter plot here]}
  \vspace{2cm}}}
  \caption{Predicted $\hat{\mu}$ vs.\ actual log-revenue on the test set (RMSE = 2.22)}
\end{figure}

\begin{figure}[H]
  \centering
  \fbox{\parbox{0.6\textwidth}{\centering\vspace{2cm}
    \textit{[Insert predictive uncertainty histogram here]}
  \vspace{2cm}}}
  \caption{Distribution of predictive $\hat{\sigma}$ per movie.
  $\hat{\sigma}$ ranges from 2.09 to 2.54, reflecting the model's
  limited ability to assign confident predictions from genre/keywords alone.}
\end{figure}

\subsection{Posterior Genre Weights}

The posterior mean $\beta$ for each genre represents its effect on
log-revenue. To aid interpretation, weights are converted to
\textbf{dollar effects relative to the training mean baseline}:

\begin{equation}
  \Delta\text{revenue}_g
  = \text{baseline} \times \bigl(\exp(\beta_g) - 1\bigr),
  \quad
  \text{baseline} = \exp\!\left(\frac{1}{N_\text{tr}}\sum_i y_i\right)
\end{equation}

where the baseline is the geometric mean revenue of the training set.

\begin{figure}[H]
  \centering
  \fbox{\parbox{0.6\textwidth}{\centering\vspace{3cm}
    \textit{[Insert genre revenue effect bar chart here]}
  \vspace{3cm}}}
  \caption{Posterior genre revenue effects relative to training mean baseline.
  Green bars = positive effect (genre earns above average),
  red bars = negative effect. Error bars show $\pm 2\sigma$ of the posterior.}
\end{figure}

\subsection{Posterior Keyword Topic Weights}

Each of the 20 SVD keyword topics has a posterior weight $\beta_{\text{kw},k}$.
The sign and magnitude indicate whether that cluster of keywords is associated
with above- or below-average revenue.

\begin{table}[H]
\centering
\small
\begin{tabular}{rlll}
\toprule
\textbf{Topic} & \textbf{$\hat{\beta}$} & \textbf{Top keywords} & \textbf{Interpretation} \\
\midrule
 1 & $+0.432$ & duringcreditsstinger, 3d, sequel    & Franchise/sequel films \\
 6 & $+0.266$ & aftercreditsstinger, musical, 3d    & Cinematic universe films \\
11 & $+0.150$ & 3d, sequel, superhero, marvel\_comic & Superhero franchise \\
 3 & $+0.194$ & sport, biography, based\_on\_novel   & Prestige biopics \\
 4 & $-0.229$ & sport, biography, independent\_film  & Low-budget sports films \\
16 & $-0.219$ & love, sex, sequel, nudity            & Adult-oriented content \\
12 & $-0.182$ & murder, based\_on\_novel, musical     & Mixed signals \\
\bottomrule
\end{tabular}
\caption{Selected keyword topic weights (full table of all 20 topics in notebook output)}
\end{table}

\begin{figure}[H]
  \centering
  \fbox{\parbox{0.6\textwidth}{\centering\vspace{3cm}
    \textit{[Insert keyword topic weight bar chart here]}
  \vspace{3cm}}}
  \caption{Posterior keyword topic weights $\hat{\beta}_\text{kw}$ with $\pm 2\sigma$ intervals.}
\end{figure}

% ─────────────────────────────────────────────────────────────────────────────
\section{Model Output and Bayesian Aggregation}

\subsection{Exported Output}

The model exports \texttt{model4\_posterior.csv} with one row per movie:

\begin{table}[H]
\centering
\begin{tabular}{ll}
\toprule
\textbf{Column} & \textbf{Meaning} \\
\midrule
\texttt{id}               & TMDB movie ID \\
\texttt{title}            & Movie title \\
\texttt{log\_revenue}     & Actual log-revenue (ground truth) \\
\texttt{mu\_genre}        & Posterior predictive mean $\hat{\mu}$ (log-revenue) \\
\texttt{sigma\_genre}     & Posterior predictive std $\hat{\sigma}$ (log-revenue) \\
\texttt{precision\_genre} & $1/\hat{\sigma}^2$ — weight for Bayesian fusion \\
\bottomrule
\end{tabular}
\caption{Columns in \texttt{model4\_posterior.csv}}
\end{table}

\subsection{Precision-Weighted Fusion}

The output feeds into a Bayesian aggregator (product-of-experts) that
combines $M$ models. Each model $m$ contributes $(\hat{\mu}_m, \hat{\sigma}_m)$.
The fused posterior is:

\begin{align}
  \sigma_\text{agg}^{-2} &= \sum_{m=1}^{M} \hat{\sigma}_m^{-2} \\[6pt]
  \mu_\text{agg} &= \sigma_\text{agg}^2 \sum_{m=1}^{M}
    \frac{\hat{\mu}_m}{\hat{\sigma}_m^2}
\end{align}

Model 4 contributes one term. A larger $\hat{\sigma}_\text{genre}$
(more uncertain) gives a smaller $1/\hat{\sigma}^2$ (lower weight),
so the aggregator automatically down-weights this model when it is
less confident.

% ─────────────────────────────────────────────────────────────────────────────
\section{Predicting Revenue for New Movies}

A \texttt{predict\_revenue()} function is provided in the notebook.
It takes genre names, keyword tags, and runtime, and returns a
full predictive distribution:

\begin{lstlisting}[language=Python, caption=Prediction function usage]
result = predict_revenue(
    genres      = ["Action", "Adventure", "Science Fiction"],
    keywords    = ["sequel", "superhero", "marvel_comic", "3d"],
    runtime_min = 150,
    release_year= 2024,
)
# Returns:
# median_2024_USD    ~ predicted revenue (50th percentile)
# p5_2024_USD        ~ lower bound (5th percentile)
# p95_2024_USD       ~ upper bound (95th percentile)
# log_rev_mu/sigma   ~ raw output for Bayesian aggregator
\end{lstlisting}

\noindent\textbf{Limitations:}
\begin{itemize}
  \item Genres outside the top 15 are ignored (mapped to zero features).
  \item Keywords not seen $\geq 3$ times in training are ignored by TF-IDF.
  \item The model does not use budget, cast, or marketing — its uncertainty
        ($\hat{\sigma} \approx 2.1$--$2.5$ in log-scale) reflects this.
  \item The prediction should be interpreted alongside other models in the
        aggregator, not as a standalone forecast.
\end{itemize}

% ─────────────────────────────────────────────────────────────────────────────
\section{Software and Reproducibility}

\begin{table}[H]
\centering
\begin{tabular}{ll}
\toprule
\textbf{Library} & \textbf{Version} \\
\midrule
Python   & 3.x \\
PyTorch  & 2.0.0 \\
Pyro-PPL & 1.9.1 \\
scikit-learn & --- \\
pandas / numpy & --- \\
\bottomrule
\end{tabular}
\caption{Software dependencies}
\end{table}

\noindent Random seed 42 is set for both PyTorch and Pyro. Results are fully
reproducible given the same data files and software versions.
The notebook is available at:
\url{https://github.com/MikkelDitlevOlsen1/Model-based-machine-learning}

\end{document}
